\documentclass[a4paper]{arrowhead}

\usepackage[yyyymmdd]{datetime}
\usepackage{etoolbox}
\usepackage[T1]{fontenc}
\usepackage[utf8]{inputenc}
\usepackage{multirow}
\usepackage{longtable}
\usepackage{amssymb}
\usepackage{float}
\usepackage{tikz}
\usetikzlibrary{positioning,arrows.meta,fit,calc,shapes.misc,shapes.geometric}

\renewcommand{\dateseparator}{-}
\setlength{\emergencystretch}{3em}
\lstset{basicstyle=\ttfamily\footnotesize,breaklines=true,columns=fullflexible,frame=single,rulecolor=\color{gray!50},backgroundcolor=\color{gray!6}}
\tikzset{
  ahsys/.style={draw=ArrowheadBlue,very thick,fill=ArrowheadBlue!6,rounded corners=2pt,minimum width=3.2cm,minimum height=1.1cm,align=center,font=\small\bfseries},
  ahext/.style={draw=ArrowheadGrey,thick,dashed,fill=gray!5,rounded corners=2pt,minimum width=3.2cm,minimum height=1.0cm,align=center,font=\small},
  ahsrv/.style={draw=ArrowheadPurple,thick,fill=ArrowheadPurple!8,rounded corners=6pt,minimum height=0.55cm,inner xsep=4pt,align=center,font=\scriptsize\ttfamily},
  ahcsm/.style={draw=ArrowheadTeal,thick,fill=ArrowheadTeal!8,rounded corners=6pt,minimum height=0.55cm,inner xsep=4pt,align=center,font=\scriptsize\ttfamily},
  ahdb/.style={draw=ArrowheadOrange,thick,fill=ArrowheadOrange!8,cylinder,shape border rotate=90,aspect=0.25,minimum height=0.9cm,align=center,font=\scriptsize},
  ahflow/.style={-{Stealth[length=2.2mm]},thick,ArrowheadGrey},
  ahlabel/.style={font=\scriptsize,fill=white,inner sep=1pt,align=center},
}

%% Special references
\newcommand{\mref}[1]{{\textcolor{ArrowheadPurple}{\hyperref[sec:model:#1]{#1}}}}
\newcommand{\oref}[1]{{\textcolor{ArrowheadBlue}{\hyperref[sec:operations:#1]{#1}}}}
\newcommand{\pdef}[1]{{\textcolor{ArrowheadGrey}{#1\label{sec:model:primitives:#1}\label{sec:model:primitives:#1s}\label{sec:model:primitives:#1es}}}}
\newcommand{\pref}[1]{{\textcolor{ArrowheadGrey}{\hyperref[sec:model:primitives:#1]{#1}}}}

\newrobustcmd\osubsection[3]{
  \addtocounter{subsection}{1}
  \addcontentsline{toc}{subsection}{\protect\numberline{\thesubsection}operation \textcolor{ArrowheadBlue}{#1}}
  \subsection*{
    \thesubsection\quad
    operation
    \textcolor{ArrowheadBlue}{#1}
    (\notblank{#2}{\mref{#2}}{})
    \notblank{#3}{: \mref{#3}}{}
  }
  \label{sec:operations:#1}
}
\newrobustcmd\msubsection[2]{
  \addtocounter{subsection}{1}
  \addcontentsline{toc}{subsection}{\protect\numberline{\thesubsection}#1 \textcolor{ArrowheadPurple}{#2}}
  \subsection*{\thesubsection\quad#1 \textcolor{ArrowheadPurple}{#2}}
  \label{sec:model:#2} \label{sec:model:#2s} \label{sec:model:#2es}
}

\begin{document}

%% Arrowhead Document Properties
\ArrowheadTitle{workflowContext}
\ArrowheadServiceID{workflowContext}
\ArrowheadType{Service Description}
\ArrowheadTypeShort{SD}
\ArrowheadVersion{5.4.0}
\ArrowheadDate{2026-11-23}
\ArrowheadAuthor{Documentation maintenance}
\ArrowheadStatus{REVIEW DRAFT}
\ArrowheadContact{Repository maintainers}
\ArrowheadFooter{\href{www.arrowhead.eu}{www.arrowhead.eu}}
\ArrowheadSetup
%%

%% Front Page
\begin{center}
  \vspace*{1cm}
  \huge{\arrowtitle}

  \vspace*{0.2cm}
  \LARGE{\arrowtype}
  \vspace*{1cm}

  \Large{Service ID: \textit{``\arrowid''}}
  \vspace*{\fill}

  

  \vspace*{1cm}
  \vspace*{\fill}

  % Front Page Abstract
  \begin{abstract}
    This document describes the \texttt{workflowContext} service of the Eclipse Arrowhead workflow support systems, which gives executors the session parameters and the results their step may read, and stores step results. It is provided by the WorkflowManager system, service version 1.0.0, documentation edition 5.4.0.
  \end{abstract}

  \vspace*{1cm}
\end{center}

\newpage
%%

%% Table of Contents
\tableofcontents
\newpage
%%

\section{Overview}
\label{sec:overview}
The \texttt{workflowContext} service holds the session context: the parameters $\rho$ fixed at the
start of a session and the results $\xi$ stored by completed steps. An executor holding a work item for step $s$ reads
$\rho$ and $\xi|dep(s)$, the results of the steps $s$ depends on, composes the request $\kappa_s(\rho, \xi|dep(s))$ and
stores the provider's answer as $\xi(s)$ before reporting success. The context replaces ad hoc data passing between
executors and keeps the data with the WorkflowManager, from where it is returned to the product.

The rest of this document is organized as follows.
In Section~\ref{sec:operations}, we describe the abstract message operations provided by the service.
In Section~\ref{sec:model}, we end the document by presenting the data types used by the mentioned operations.

\subsection{Significant Prior Art}
The session-context service of the Arrowhead workflow kits (\texttt{github.com/palvarga/arrowhead-workflows})
and Lemma~2 of~\cite{varga2026}: for a well-formed recipe a step never runs on missing data.

\subsection{How This Service Is Meant to Be Used}
Before invoking the provider of step $s$, the executor calls \oref{get-parameters} and \oref{get-results}
for its work item; after a successful invocation it calls \oref{put-result} and only then reports \texttt{SUCCESS}
through \texttt{workflowWorkItem}.

\subsection{Important Delimitations and Dependencies}
\label{sec:delimitations}
\begin{itemize}
\item Only the holder of the work item (or a sysop) may read or write its context (403 otherwise).
\item Results are restricted to $dep(s)$; there is no way to read other results of the session.
\item A dependency skipped by a guard (extension x-guard) is returned as null. A missing result of a completed dependency would violate Lemma~2 and is refused with 409.
\item \oref{put-result} stores a pending result. When the following \texttt{SUCCESS} report also carries a result, the reported one is used.
\end{itemize}

\paragraph{Authorization.} Each operation is authorized separately; the holder check is made in addition.

\newpage
\section{Service Interface}
\label{sec:operations}
This section describes the interface of the \texttt{workflowContext} service. Each subsection names an abstract operation,
an input type and an output type, in that order. The input type is named inside parentheses, while the output type is
preceded by a colon. Input and output types are only denoted when accepted or returned, respectively, by the operation
in question. All abstract data types named in this section are defined in Section~\ref{sec:model}. The concrete
interface (HTTP/TLS/JSON) is specified in the IDD of the same name.

\begin{figure}[ht!]
  \centering
  \resizebox{0.92\textwidth}{!}{
\begin{tikzpicture}
\node[ahsys] (prov) at (4.6,0) {WorkflowManager\\[-1pt]{\normalfont\scriptsize $\ll$provider$\gg$}};
\node[ahsrv] (srv) at (1.6,0) {workflowContext};
\draw[thick,ArrowheadPurple] (srv.east) -- (prov.west);
\node[ahext] (c0) at (-3.0,0.00) {WorkflowExecutor\\[-1pt]{\scriptsize $\ll$consumer$\gg$}};
\draw[ahflow] (c0.east) -- (srv.west);
\end{tikzpicture}
  }
  \caption{SysML block view of the \texttt{workflowContext} service: provider and typical consumers.}
  \label{fig:SD}
\end{figure}

The following interface operations are available.

{\small
\begin{longtable}{|p{4.4cm}|p{5.6cm}|p{5.8cm}|}

\hline
\rowcolor{gray!33} Operation & Input & Output \\ \hline
\endhead
\oref{get-parameters} & \mref{WorkItemRef} & \mref{Parameters} \\ \hline
\oref{get-results} & \mref{WorkItemRef} & \mref{StepResults} \\ \hline
\oref{put-result} & \mref{PutResultRequest} & -- \\ \hline
\end{longtable}
}

\osubsection{get-parameters}{WorkItemRef}{Parameters}
Returns $\rho$ of the work item's session.
\osubsection{get-results}{WorkItemRef}{StepResults}
Returns $\xi|dep(s)$ for the work item's step.
\osubsection{put-result}{PutResultRequest}{}
Stores $\xi(s)$ for the work item's step.

\newpage
\section{Information Model}
\label{sec:model}
Here, all data objects that can be part of the \texttt{workflowContext} service are listed in alphabetic order. Each
subsection, which describes one type of object, begins with the \textit{struct} keyword (a collection of named fields,
each with its own data type), the \textit{enum} keyword (a closed set of values) or the \textit{alias} keyword. As a
complement to the explicitly defined types in this section, there is also a list of implicit primitive types in
Section~\ref{sec:model:primitives}.

\msubsection{alias}{Parameters}
The session parameters $\rho$.
 It is an alias of \pref{JsonObject}.
\msubsection{struct}{PutResultRequest}
Stores $\xi(s)$ for a held work item.

{\small
\begin{longtable}{|p{3.9cm}|p{4.4cm}|p{7.5cm}|}

\hline
\rowcolor{gray!33} Field & Type & Description \\ \hline
\endhead
\texttt{workItemId} & \pref{WorkItemId} & The work item. \\ \hline
\texttt{result} & \pref{JsonValue} & The step's result. \\ \hline
\end{longtable}
}

\msubsection{alias}{StepResults}
$\xi|dep(s)$: the stored results of the steps in $dep(s)$, by step name.
 It is an alias of \pref{Map}$<$\pref{Name}, \pref{JsonValue}$>$.
\msubsection{struct}{WorkItemRef}
Identifies a work item.

{\small
\begin{longtable}{|p{3.9cm}|p{4.4cm}|p{7.5cm}|}

\hline
\rowcolor{gray!33} Field & Type & Description \\ \hline
\endhead
\texttt{workItemId} & \pref{WorkItemId} & The work item. \\ \hline
\end{longtable}
}


\subsection{Primitives}
\label{sec:model:primitives}
Types and structures mentioned throughout this document that are assumed to be available to implementations of this
service. The concrete interpretation of each of these types must be provided by any IDD document claiming to implement
this service.

{\small
\begin{longtable}{|p{3.2cm}|p{12.6cm}|}

\hline
\rowcolor{gray!33} Type & Description \\ \hline
\endhead
\pdef{JsonObject} & A set of named values. \\ \hline
\pdef{JsonValue} & Any value: object, array, string, number, Boolean or null. \\ \hline
\pdef{Map} & A collection of key/value pairs of the types given in angle brackets. \\ \hline
\pdef{Name} & An identifier of a recipe, action, step, parameter or resource: a letter followed by at most 62 letters, digits, \texttt{-} or \texttt{\_}. \\ \hline
\pdef{WorkItemId} & The identifier of one attempt to run one step of one session. \\ \hline
\end{longtable}
}


\newpage
\bibliographystyle{IEEEtran}
\bibliography{bibliography}

\newpage

\section{Revision History}
\subsection{Amendments}

\noindent\begin{tabularx}{\textwidth}{| p{1cm} | p{2.2cm} | p{1.6cm} | X | p{3.6cm} |} \hline
\rowcolor{gray!33} No. & Date & Version & Subject of Amendments & Author \\ \hline
1 & 2026-11-23 & \arrowversion & Initial edition: service description of workflowContext, written from the Java reference implementation (ah5-Workflow-java). & Documentation maintenance \\ \hline
\end{tabularx}

\subsection{Quality Assurance}

\noindent\begin{tabularx}{\textwidth}{| p{1cm} | p{2.2cm} | p{1.6cm} | X |} \hline
\rowcolor{gray!33} No. & Date & Version & Approved by \\ \hline
1 & & \arrowversion & Pending review \\ \hline
\end{tabularx}

\end{document}
